\section{引言：重新思考面向 Agent 的软件设计}

随着 Model Context Protocol (MCP) 等标准的出现，LLM Agent 可以通过数百个工具来完成现实世界的任务。但一个核心问题是：\textbf{如何让这些工具发挥最大效用？}

本文是 Anthropic 工程团队在实际开发中积累的工具设计最佳实践，涵盖了从原型搭建、评估体系到工具优化的全流程，最终提炼出五大核心原则。文章的一大亮点在于：\textbf{你可以让 Agent 自己来优化它所使用的工具}。

\begin{importantbox}{核心观点}
Agent 的能力上限取决于我们赋予它的工具。为 Agent 设计工具与为人类开发者设计 API 有本质区别——我们需要从确定性系统（deterministic systems）转向为非确定性系统（non-deterministic agents）设计软件接口。
\end{importantbox}

\subsection{Tool 的本质：确定性与非确定性的桥梁}

在传统软件工程中，我们编写的是\textbf{确定性系统之间的契约}。例如，调用 \texttt{getWeather("NYC")} 每次都会以完全相同的方式获取纽约的天气数据。

但 Tool 是一种\textbf{新型软件}，它反映的是确定性系统与非确定性 Agent 之间的契约。当用户问"今天需要带伞吗？"时，Agent 可能会调用天气工具，也可能从通用知识中回答，甚至可能先追问用户所在的城市。偶尔，Agent 还可能产生幻觉（hallucination）或根本不理解如何使用某个工具。

\begin{knowledgebox}{确定性 vs 非确定性系统}
\begin{itemize}[leftmargin=1.5em]
    \item \textbf{确定性系统}：相同输入总是产生相同输出。传统函数调用、API 请求都属于此类。
    \item \textbf{非确定性系统}：即使输入相同，也可能产生不同的输出。LLM Agent 就是典型代表。
    \item \textbf{Tool 的角色}：在两者之间搭建桥梁，让 Agent 能够与确定性的外部世界交互。
\end{itemize}
\end{knowledgebox}

这意味着我们需要根本性地重新思考软件编写方式：不再像为其他开发者或系统编写函数和 API 那样编写工具和 MCP Server，而要\textbf{专门为 Agent 而设计}。

文章的核心目标是：扩大 Agent 能够有效解决问题的覆盖面，让 Agent 通过工具追求多种成功策略。一个有趣的发现是：对 Agent 最"符合人体工学"（ergonomic）的工具，往往对人类来说也出奇地直观。

\subsection{本章小结}

Tool 是连接确定性系统和非确定性 Agent 的新型软件接口。为 Agent 设计工具需要从传统 API 设计思维中跳脱出来，核心在于理解 Agent 的认知特性和行为模式，而非简单地暴露底层功能。


